\section*{Chapter \ref{ch:ll:multi-socket-machines-and-interconnects} --- Multi-Socket Machines and Interconnects}

\begin{solution}{exo:ll:multi-socket-machines-and-interconnects:1}
$16 \times 8 \times 128/130/8 = 15.75$~GB/s in each direction; two 25~Gb/s ports need 6.25~GB/s: yes, with room to spare.
\end{solution}

\begin{solution}{exo:ll:multi-socket-machines-and-interconnects:2}
$\lceil 300/64 \rceil = 5$ lines if the packet starts on a line boundary, 6 if it does not.
\end{solution}

\begin{solution}{exo:ll:multi-socket-machines-and-interconnects:3}
The card is on node 1 and the strategy on node 0: remote. $(3 + 1)(138 - 59) = \qty{316}{\nano\second}$ per message.
\end{solution}

\begin{solution}{exo:ll:multi-socket-machines-and-interconnects:4}
First touch put the ring's pages on node 0, so every write from node 1 goes across the interconnect to node 0's memory or
caches, and the reader pays again if it is on node 1. Allocate and zero the ring in the feed thread after pinning it, or
bind the ring's memory to node 1.
\end{solution}

\begin{solution}{exo:ll:multi-socket-machines-and-interconnects:5}
Two neighbouring virtual CPUs, for example 0 and 1: a hand-off of about 40 nanoseconds one way on the committed
matrix, against a median of about 190 for other pairs, four to five times more. The choice holds only as long as the
hypervisor's mapping does.
\end{solution}

\begin{solution}{exo:ll:multi-socket-machines-and-interconnects:6}
The sibling shares the core's execution units, first-level caches and predictors, so any work on it slows and \omterm{def:ll:where-latency-comes-from:tail}{jitters} the
hot thread. On the fixture the three hot threads leave CPUs 29, 30 and 31 idle: three of node 1's sixteen CPUs.
\end{solution}

\begin{solution}{exo:ll:multi-socket-machines-and-interconnects:7}
Share one core between the two threads that exchange the least data per message and least often compete for it: the
gateway and the book builder are poor candidates (both on every message); the pair that exchanges data only through
another thread, for example the book builder and the gateway when the strategy sits between them, is the least bad.
\texttt{check} must accept a declared pair of roles on one physical core, and nothing else.
\end{solution}

\begin{solution}{exo:ll:multi-socket-machines-and-interconnects:8}
\texttt{taskset} pins the thread but says nothing about the card: nobody checked that CPU 5 is on the card's node, the
choice breaks when the server or the slot changes, the process's memory may have been first touched elsewhere, and
threads created later inherit the mask rather than getting their own cores. Derive the placement from the topology and
check it at start-up.
\end{solution}

\begin{solution}{pb:ll:multi-socket-machines-and-interconnects:1}
\begin{enumerate}
\item \path{/sys/bus/pci/devices/<address>/numa_node} for the card (found from
\path{/sys/class/net/<name>/device}), and \path{/sys/devices/system/node/node*/cpulist} for the CPUs.
\item CPUs 16 to 31; CPU 17's siblings are 17 and 25.
\item Feed on 21, strategy on 22, gateway on 23; CPUs 29, 30 and 31 idle.
\item The kernel, timers and many drivers favour CPU 0; keeping its core for housekeeping keeps that work off the hot
path.
\item $(2 + 2)(138 - 59) = \qty{316}{\nano\second}$.
\item $316/2\,500 = 12.6\%$.
\item Each book access now misses to node 0's memory across the interconnect: the threads are local to the card but their
data is remote.
\item The book: its accesses are many per message and dependent (a tree or ladder walk), while the card's lines are few.
\item One line each way: $2 \times 59 = \qty{118}{\nano\second}$ within a socket, $2 \times 138 = \qty{276}{\nano\second}$
across.
\item About 40 nanoseconds one way between neighbouring CPUs and about 190 between the others, on the committed
measurement.
\item Choose the pair with the lowest measured hand-off, among CPUs on the card's node, and verify it after every boot.
\item The hypervisor maps virtual CPUs to physical cores as it sees fit, and the mapping can differ between boots or change
while running.
\item \textbf{Named result.} The wrong socket adds about \qty{316}{\nano\second} per message, 12.6\% of a
\qty{2.5}{\micro\second} wire-to-wire budget, before counting data first touched on the wrong node.
\item Read the topology; compute the plan; pin each hot thread before it allocates; allocate and touch its memory from it;
move other devices' interrupts off the hot cores; check the plan at start-up and refuse to trade if it fails.
\item Read each thread's CPU (\path{/proc/<pid>/task/*/stat} or \texttt{sched\_getcpu} logged by the thread) and each
region's node (\path{/proc/<pid>/numa_maps}) and compare with the plan.
\item The plan is recomputed from the topology and the process moves with the card; a start-up check catches it if nobody
does.
\item Those of the disks, other network cards, and timers; the card's own interrupts, if it uses them, go to the node's
housekeeping cores.
\item The firmware did not report the card's node: the plan falls back to a default and must say so, and a human must
confirm.
\item The placement plan (the topology, each thread's CPU and node), with the machine and software versions.
\item Because where a thread and its data sit decides how many interconnect crossings each message pays.
\end{enumerate}
\end{solution}

\begin{solution}{iq:ll:multi-socket-machines-and-interconnects:1}
Each socket has its own memory; memory attached to another socket is slower to reach over the interconnect. A
latency-sensitive process must keep its hot threads, their memory and the network card on one node.
\iqlookfor{local versus remote memory and the placement consequence.}
\end{solution}

\begin{solution}{iq:ll:multi-socket-machines-and-interconnects:2}
Linux places a page on the node of the CPU that first writes it; if the main thread initialises the data at start-up and
the worker runs on another node, the worker's data is remote. Pin first, then touch, or bind explicitly.
\iqlookfor{that initialisation, not allocation, decides placement.}
\end{solution}

\begin{solution}{iq:ll:multi-socket-machines-and-interconnects:3}
Find its PCI address from \path{/sys/class/net/<if>/device}, then read \texttt{numa\_node} and \texttt{local\_cpulist}
under \path{/sys/bus/pci/devices/<address>/}; \texttt{lstopo} shows it graphically.
\iqlookfor{sysfs, and awareness that $-1$ means unknown.}
\end{solution}

\begin{solution}{iq:ll:multi-socket-machines-and-interconnects:4}
On the same socket, ideally on cores that share a cache level, each on its own physical core with siblings idle, on the
network card's node: a hand-off across sockets costs more than twice as much.
\iqlookfor{core-to-core latency as the cost of a hand-off.}
\end{solution}

\begin{solution}{iq:ll:multi-socket-machines-and-interconnects:5}
The card writes packets into host memory by itself, without the processor copying them; with \omterm{def:ll:multi-socket-machines-and-interconnects:dca}{direct cache access} they
land in the last-level cache of the card's socket, so the core reading them hits in cache instead of memory, provided the
core is on that socket.
\iqlookfor{device-driven transfers and their locality.}
\end{solution}

\begin{solution}{iq:ll:multi-socket-machines-and-interconnects:6}
At start-up, read the topology; derive each thread's CPU from the card's node by a fixed rule (one physical core each,
siblings idle, CPU 0 excluded); pin each thread before it allocates; allocate and touch its memory from it or bind it;
verify the placement from \path{/proc} and refuse to trade on a mismatch; log the plan with every measurement.
\iqlookfor{derived rather than hand-typed placement, and a check that refuses to start.}
\end{solution}
